% ===== Heap / Priority Queue (7 problems) =====
\section{Heap / Priority Queue \texorpdfstring{{\normalsize\color{muted}$\cdot$ 7 problems}}{- 7 problems}}
\markboth{Heap / Priority Queue}{}

\begin{problem}{703}{Kth Largest Element in a Stream}{Easy}
\summary{Class that returns the k-th largest after each added value.}
\keyidea{Keep a min-heap of size k; its top is the k-th largest.}
\complexity{O(log k) per add, O(k) space}
\begin{code}
class KthLargest:
    def __init__(self, k: int, nums: List[int]):
        self.k, self.heap = k, nums
        heapq.heapify(self.heap)
        while len(self.heap) > k:
            heapq.heappop(self.heap)

    def add(self, val: int) -> int:
        heapq.heappush(self.heap, val)
        if len(self.heap) > self.k:
            heapq.heappop(self.heap)
        return self.heap[0]
\end{code}
\end{problem}

\begin{problem}{1046}{Last Stone Weight}{Easy}
\summary{Repeatedly smash the two heaviest stones; return the last stone's weight (or 0).}
\keyidea{Max-heap via negated values.}
\complexity{O(n log n) time, O(n) space}
\begin{code}
class Solution:
    def lastStoneWeight(self, stones: List[int]) -> int:
        h = [-s for s in stones]
        heapq.heapify(h)
        while len(h) > 1:
            a, b = -heapq.heappop(h), -heapq.heappop(h)
            if a != b:
                heapq.heappush(h, -(a - b))
        return -h[0] if h else 0
\end{code}
\end{problem}

\begin{problem}{973}{K Closest Points to Origin}{Medium}
\summary{Return the k points nearest to (0, 0).}
\keyidea{Max-heap of size k on squared distance (negated); drop the farthest when over k.}
\complexity{O(n log k) time, O(k) space}
\begin{code}
class Solution:
    def kClosest(self, points: List[List[int]], k: int) -> List[List[int]]:
        h = []
        for x, y in points:
            heapq.heappush(h, (-(x * x + y * y), x, y))
            if len(h) > k:
                heapq.heappop(h)
        return [[x, y] for _, x, y in h]
\end{code}
\end{problem}

\begin{problem}{215}{Kth Largest Element in an Array}{Medium}
\summary{Return the k-th largest value in an unsorted array.}
\keyidea{Min-heap of size k (or Quickselect for average O(n)).}
\complexity{O(n log k) time, O(k) space}
\begin{code}
class Solution:
    def findKthLargest(self, nums: List[int], k: int) -> int:
        h = []
        for n in nums:
            heapq.heappush(h, n)
            if len(h) > k:
                heapq.heappop(h)
        return h[0]
\end{code}
\end{problem}

\begin{problem}{621}{Task Scheduler}{Medium}
\summary{Same tasks need n idle slots between them; minimum total time units to run all.}
\keyidea{Formula: (maxCount - 1) * (n + 1) + numberOfTasksWithMaxCount, but at least len(tasks).}
\complexity{O(N) time, O(26) space}
\begin{code}
class Solution:
    def leastInterval(self, tasks: List[str], n: int) -> int:
        count = Counter(tasks)
        mx = max(count.values())
        num_mx = sum(1 for v in count.values() if v == mx)
        return max(len(tasks), (mx - 1) * (n + 1) + num_mx)
\end{code}
\end{problem}

\begin{problem}{355}{Design Twitter}{Medium}
\summary{postTweet, follow, unfollow, getNewsFeed (10 most recent from self + followees).}
\keyidea{Global decreasing time counter; merge the followees' tweet lists with a heap (k-way merge).}
\complexity{getNewsFeed O(F log F) for F followees; O(users + tweets) space}
\begin{code}
class Twitter:
    def __init__(self):
        self.time = 0
        self.tweets = defaultdict(list)    # user -> [(time, tweetId)]
        self.following = defaultdict(set)

    def postTweet(self, userId: int, tweetId: int) -> None:
        self.tweets[userId].append((self.time, tweetId))
        self.time -= 1                     # smaller = newer (min-heap)

    def getNewsFeed(self, userId: int) -> List[int]:
        res, h = [], []
        users = self.following[userId] | {userId}
        for u in users:
            if self.tweets[u]:
                i = len(self.tweets[u]) - 1
                t, tid = self.tweets[u][i]
                h.append((t, tid, u, i - 1))
        heapq.heapify(h)
        while h and len(res) < 10:
            t, tid, u, i = heapq.heappop(h)
            res.append(tid)
            if i >= 0:
                t2, tid2 = self.tweets[u][i]
                heapq.heappush(h, (t2, tid2, u, i - 1))
        return res

    def follow(self, followerId: int, followeeId: int) -> None:
        if followerId != followeeId:
            self.following[followerId].add(followeeId)

    def unfollow(self, followerId: int, followeeId: int) -> None:
        self.following[followerId].discard(followeeId)
\end{code}
\end{problem}

\begin{problem}{295}{Find Median from Data Stream}{Hard}
\summary{Support addNum and findMedian on a growing stream.}
\keyidea{Max-heap for the lower half, min-heap for the upper half; keep sizes within 1.}
\complexity{addNum O(log n), findMedian O(1)}
\begin{code}
class MedianFinder:
    def __init__(self):
        self.small, self.large = [], []    # small = max-heap (negated)

    def addNum(self, num: int) -> None:
        heapq.heappush(self.small, -num)
        heapq.heappush(self.large, -heapq.heappop(self.small))
        if len(self.large) > len(self.small):
            heapq.heappush(self.small, -heapq.heappop(self.large))

    def findMedian(self) -> float:
        if len(self.small) > len(self.large):
            return -self.small[0]
        return (-self.small[0] + self.large[0]) / 2
\end{code}
\end{problem}
